% GENERATED FILE. Edit canonical lessons, then run npm run generate:books.
% canonical-source-hash: fnv1a64:05dccc125f65c262
% canonical-lessons: SA-C186-komalah, SA-C186-vamanah, SA-C186-visalah, SA-C186-suksmah, SA-C186-sthirah

\chapter{Tender, Dwarfish, Vast, Fine, Steady}
\label{ch:sa-tender-dwarfish-vast-fine-steady}
\hlchaptermodality{186}

\begin{chapteropening}
\textbf{By the end of this chapter:} \emph{I can say tender, dwarfish, vast, fine and steady.}

Complete the last lesson of chapter 186: I can say tender, dwarfish, vast, fine and steady.
\end{chapteropening}

\section[\texorpdfstring{\sk{कोमलः} (komalaḥ)}{कोमलः (komalaḥ)}]{\sk{कोमलः} (komalaḥ) — tender, soft}

\label{lesson:SA-C186-komalah}



\begin{quote}
Before the new one: say the Sanskrit for surely, then the Sanskrit for certainly.
\end{quote}

\subsection*{You'll want to know: \sk{कोमलः}}
\textbf{\sk{कोमलः}} — \emph{komalaḥ} — \textquotedblleft{}tender, soft\textquotedblright{}.

\begin{grammarlens}[title={What you've built}]
One more word for this chapter.
\end{grammarlens}

\subsection*{Your turn}
\begin{itemize}
\raggedright
  \item \emph{Say it:} \emph{komalaḥ}
  \item \emph{Say it:} \emph{komalaḥ}, once more
  \item \emph{Say it:} say \emph{komalaḥ}
  \item \emph{Recall:} say the Sanskrit for surely, then the Sanskrit for certainly, then say \emph{komalaḥ} again
\end{itemize}

\begin{tcolorbox}[breakable,colback=teal!4,colframe=teal!35!black,title={Before you move on}]
What does \sk{कोमलः} mean? (\textquotedblleft{}Tender, soft\textquotedblright{}.) Say it once more.
\end{tcolorbox}
\section[\texorpdfstring{\sk{वामनः} (vāmanaḥ)}{वामनः (vāmanaḥ)}]{\sk{वामनः} (vāmanaḥ) — dwarfish, short}

\label{lesson:SA-C186-vamanah}



\begin{quote}
Before the new one: say the Sanskrit for certainly, then the Sanskrit for tender.
\end{quote}

\subsection*{You'll want to know: \sk{वामनः}}
\textbf{\sk{वामनः}} — \emph{vāmanaḥ} — \textquotedblleft{}dwarfish, short\textquotedblright{}.

\begin{grammarlens}[title={What you've built}]
One more word for this chapter.
\end{grammarlens}

\subsection*{Your turn}
\begin{itemize}
\raggedright
  \item \emph{Say it:} \emph{vāmanaḥ}
  \item \emph{Say it:} \emph{vāmanaḥ}, once more
  \item \emph{Say it:} say \emph{vāmanaḥ}
  \item \emph{Recall:} say the Sanskrit for certainly, then the Sanskrit for tender, then say \emph{vāmanaḥ} again
\end{itemize}

\begin{tcolorbox}[breakable,colback=teal!4,colframe=teal!35!black,title={Before you move on}]
What does \sk{वामनः} mean? (\textquotedblleft{}Dwarfish, short\textquotedblright{}.) Say it once more.
\end{tcolorbox}
\section[\texorpdfstring{\sk{विशालः} (viśālaḥ)}{विशालः (viśālaḥ)}]{\sk{विशालः} (viśālaḥ) — vast, broad}

\label{lesson:SA-C186-visalah}



\begin{quote}
Before the new one: say the Sanskrit for tender, then the Sanskrit for dwarfish.
\end{quote}

\subsection*{You'll want to know: \sk{विशालः}}
\textbf{\sk{विशालः}} — \emph{viśālaḥ} — \textquotedblleft{}vast, broad\textquotedblright{}.

\begin{grammarlens}[title={What you've built}]
One more word for this chapter.
\end{grammarlens}

\subsection*{Your turn}
\begin{itemize}
\raggedright
  \item \emph{Say it:} \emph{viśālaḥ}
  \item \emph{Say it:} \emph{viśālaḥ}, once more
  \item \emph{Say it:} say \emph{viśālaḥ}
  \item \emph{Recall:} say the Sanskrit for tender, then the Sanskrit for dwarfish, then say \emph{viśālaḥ} again
\end{itemize}

\begin{tcolorbox}[breakable,colback=teal!4,colframe=teal!35!black,title={Before you move on}]
What does \sk{विशालः} mean? (\textquotedblleft{}Vast, broad\textquotedblright{}.) Say it once more.
\end{tcolorbox}
\section[\texorpdfstring{\sk{सूक्ष्मः} (sūkṣmaḥ)}{सूक्ष्मः (sūkṣmaḥ)}]{\sk{सूक्ष्मः} (sūkṣmaḥ) — fine, subtle}

\label{lesson:SA-C186-suksmah}



\begin{quote}
Before the new one: say the Sanskrit for dwarfish, then the Sanskrit for vast.
\end{quote}

\subsection*{You'll want to know: \sk{सूक्ष्मः}}
\textbf{\sk{सूक्ष्मः}} — \emph{sūkṣmaḥ} — \textquotedblleft{}fine, subtle\textquotedblright{}.

\begin{grammarlens}[title={What you've built}]
One more word for this chapter.
\end{grammarlens}

\subsection*{Your turn}
\begin{itemize}
\raggedright
  \item \emph{Say it:} \emph{sūkṣmaḥ}
  \item \emph{Say it:} \emph{sūkṣmaḥ}, once more
  \item \emph{Say it:} say \emph{sūkṣmaḥ}
  \item \emph{Recall:} say the Sanskrit for dwarfish, then the Sanskrit for vast, then say \emph{sūkṣmaḥ} again
\end{itemize}

\begin{tcolorbox}[breakable,colback=teal!4,colframe=teal!35!black,title={Before you move on}]
What does \sk{सूक्ष्मः} mean? (\textquotedblleft{}Fine, subtle\textquotedblright{}.) Say it once more.
\end{tcolorbox}
\section[\texorpdfstring{\sk{स्थिरः} (sthiraḥ)}{स्थिरः (sthiraḥ)}]{\sk{स्थिरः} (sthiraḥ) — steady}

\label{lesson:SA-C186-sthirah}



\begin{quote}
Before the new one: say the Sanskrit for vast, then the Sanskrit for fine.
\end{quote}

\subsection*{You'll want to know: \sk{स्थिरः}}
\textbf{\sk{स्थिरः}} — \emph{sthiraḥ} — \textquotedblleft{}steady\textquotedblright{}.

\begin{grammarlens}[title={What you've built}]
One more word for this chapter.
\end{grammarlens}

\subsection*{Your turn}
\begin{itemize}
\raggedright
  \item \emph{Say it:} \emph{sthiraḥ}
  \item \emph{Say it:} \emph{sthiraḥ}, once more
  \item \emph{Say it:} say \emph{sthiraḥ}
  \item \emph{Recall:} say the Sanskrit for vast, then the Sanskrit for fine, then say \emph{sthiraḥ} again
\end{itemize}

\begin{tcolorbox}[breakable,colback=teal!4,colframe=teal!35!black,title={Before you move on}]
What does \sk{स्थिरः} mean? (\textquotedblleft{}Steady\textquotedblright{}.) Say it once more.
\end{tcolorbox}
